\chapter{Introdução}

O consumo de energia elétrica muda de uma casa para outra. Ele depende dos
moradores, do clima, do tamanho da casa e dos aparelhos.

Este trabalho utiliza aprendizado supervisionado \cite{ia2013}, no qual o algoritmo recebe
pares de entrada e saída para aprender uma função $h$, que estima o consumo de energia para entradas novas.
Aqui a saída é o consumo em kWh, um valor contínuo. Por isso, o
problema é de regressão.

O objetivo é prever o consumo anual de cada casa com regressão linear
múltipla. Nela, $h$ soma as entradas multiplicadas por pesos, e os pesos são
escolhidos para deixar o erro de previsão o menor possível. Os dados vêm do
RECS 2020, uma pesquisa da U.S. Energy Information Administration com casas
dos Estados Unidos \cite{eia2020recs}.

O Capítulo~2 mostra os dados e o modelo. O Capítulo~3 mostra os resultados, e
o Capítulo~4, a conclusão.
